\chapter{Costruire un programma}
In un qualsiasi programma, ci sono alcune cose da definire e considerare. Infatti, possiamo a tutti gli effetti considerarlo un sistema.

\section{Dimensione e Complessità}
Si possono definire due grandi problemi per la dimensione e la complessità:
\begin{enumerate}
    \item \textbf{\textcolor{brown}{Breadth issues} (problemi di ampiezza):}
    \begin{itemize}
        \item \textbf{Funzioni principali}, generiche di cosa voglio che faccia.
        \item \textbf{Features} all'interno di ogni funzione.
        \item \textbf{Interfacce} verso sistemi esterni.
        \item \textbf{Utenti} connessi simultaneamente.
        \item Tipi di \textbf{dati e strutture} in cui vengono raccolti.
    \end{itemize}
    \item \textbf{\textcolor{brown}{Depth issues} (problemi di profondità):}
    \begin{itemize}
        \item \textbf{Collegamenti e relazioni} tra gli oggetti: condivisione dei dati e di controllo, gerarchie, sequenziale, iterativo, ricorsivo.
    \end{itemize}
\end{enumerate}

\subsection{Divide et Impera} 
Si deve suddividere il problema in problemi minori più facilmente risolvibili.
    \begin{itemize}
        \item \textbf{Decomposizione.}
        \item \textbf{Modularizzazione.}
        \item \textbf{Separazione}
        \item \textbf{Iterazioni} incrementali.
    \end{itemize}

Quando si lavora in team sullo sviluppo di un progetto complesso bisogna mettersi d'accordo su vari aspetti:
\begin{itemize}
    \item Linguaggio di programmazione, tecnologie e ambiente di sviluppo.
    \item Sistema di controllo di versione, suddivisione e implementazione di task anche collegati.
\end{itemize}

\section{CI/CD (Continuous Integration/Continuous Deployment}
Ogni task una volta completata viene integrata e testato tutto.
Alcuni progetti diventano un fallimento perché è difficile stimare l'\textbf{\textcolor{blue}{effort}} reale per sviluppare tutte le \textbf{features} richieste. Quando si passa all'integrazione di una nuova \textbf{features} è necessaria la comunicazione tra il team e a volte può fallire.
Quando si gestisce un sistema ampio è necessaria la coordinazione tra le \textcolor{red}{3P} perché c'è un importante \textbf{effort} di coordinazione:
\begin{enumerate}
    \item \textbf{\textcolor{red}{Processi}} e metodologie utilizzate: migliorare la qualità del software e la produttività degli sviluppatori.
    \item \textbf{\textcolor{red}{Prodotto}} finale e artefatti intermedi: al giorno d'oggi non basta consegnare l'eseguibile finale ma anche documentazioni e altro.
    \item \textbf{\textcolor{red}{Persone}} (sviluppatori, personale di supporto e utenti): coordinazione di queste ultime e attento studio sulle skill di ciascuno.
\end{enumerate}


